
\FigureLayoutDeclare{img/2016/national_paper_2_liberal/q19_fig1.png}{14.0cm}{56bp 61bp 101bp 69bp}

\chapter{2016年全国II卷（文）}

\section{单选题}

\begin{problem}
已知集合 \(A=\{1,2,3\}\)，\(B=\{x\mid x^2<9\}\)，则 \(A\cap B=\)（\quad）

\choices
    {\(\{-2,-1,0,1,2,3\}\)}
    {\(\{-2,-1,0,1,2\}\)}
    {\(\{1,2,3\}\)}
    {\(\{1,2\}\)}
\end{problem}

\begin{answer}
D
\end{answer}

\begin{problem}
设复数 \(z\) 满足 \(z+\i=3-\i\)，则 \(\overline z=\)（\quad）

\choices
    {\(-1+2\i\)}
    {\(1-2\i\)}
    {\(3+2\i\)}
    {\(3-2\i\)}
\end{problem}

\begin{answer}
C
\end{answer}


\begin{problem}
函数 \(y=A\sin(\omega x+\varphi)\) 的部分图象如图所示，则（\quad）

\begingroup
\bitmapfigure[max width=0.24\linewidth]{img/2016/national_paper_2_liberal/q3_fig1.png}
\endgroup

\choices
    {\(y=2\sin\left(2x-\dfrac{\pi}{6}\right)\)}
    {\(y=2\sin\left(2x-\dfrac{\pi}{3}\right)\)}
    {\(y=2\sin\left(x+\dfrac{\pi}{6}\right)\)}
    {\(y=2\sin\left(x+\dfrac{\pi}{3}\right)\)}
\end{problem}

\begin{answer}
A
\end{answer}


\begin{problem}
体积为 \(8\) 的正方体的顶点都在同一球面上，则该球面的表面积为（\quad）

\choices
    {\(12\pi\)}
    {\(\dfrac{32}{3}\pi\)}
    {\(8\pi\)}
    {\(4\pi\)}
\end{problem}

\begin{answer}
A
\end{answer}


\begin{problem}
设 \(F\) 为抛物线 \(C:y^2=4x\) 的焦点，曲线 \(y=\dfrac{k}{x}\)（\(k>0\)）与 \(C\) 交于点 \(P\)，且 \(PF\perp x\) 轴，则 \(k=\)（\quad）

\choices
    {\(\dfrac12\)}
    {\(1\)}
    {\(\dfrac32\)}
    {\(2\)}
\end{problem}

\begin{answer}
D
\end{answer}


\begin{problem}
圆 \(x^2+y^2-2x-8y+13=0\) 的圆心到直线 \(ax+y-1=0\) 的距离为 \(1\)，则 \(a=\)（\quad）

\choices
    {\(-\dfrac43\)}
    {\(-\dfrac34\)}
    {\(\sqrt3\)}
    {\(2\)}
\end{problem}

\begin{answer}
A
\end{answer}


\begin{problem}
如图是由圆柱与圆锥组合而成的几何体的三视图，则该几何体的表面积为（\quad）

\begingroup
\bitmapfigure[float=inline,max width=0.42\linewidth]{img/2016/national_paper_2_liberal/q7_fig1.png}
\endgroup

\choices
    {\(20\pi\)}
    {\(24\pi\)}
    {\(28\pi\)}
    {\(32\pi\)}
\end{problem}

\begin{answer}
C
\end{answer}


\begin{problem}
某路口人行横道的信号灯为红灯和绿灯交替出现，红灯持续时间为 \(40\text{ 秒}\). 若一名行人来到该路口遇到红灯，则至少需要等待 \(15\text{ 秒}\) 才出现绿灯的概率为（\quad）

\choices
    {\(\dfrac{7}{10}\)}
    {\(\dfrac{5}{8}\)}
    {\(\dfrac{3}{8}\)}
    {\(\dfrac{3}{10}\)}
\end{problem}

\begin{answer}
B
\end{answer}


\begin{problem}
中国古代有计算多项式值的秦九韶算法，如图是实现该算法的程序框图. 执行该程序框图，若输入的 \(x=2\)，\(n=2\)，依次输入的 \(a\) 为 \(2\)，\(2\)，\(5\)，则输出的 \(s=\)（\quad）

\begingroup
\texfigure[max width=6.0cm]{img_repaint/2016/national_paper_2_liberal/q9_fig1.tex}
\endgroup

\choices
    {\(7\)}
    {\(12\)}
    {\(17\)}
    {\(34\)}
\end{problem}

\begin{answer}
C
\end{answer}


\begin{problem}
下列函数中，其定义域和值域分别与函数 \(y=10^{\lg x}\) 的定义域和值域相同的是（\quad）

\choices
    {\(y=x\)}
    {\(y=\lg x\)}
    {\(y=2^x\)}
    {\(y=\dfrac{1}{\sqrt{x}}\)}
\end{problem}

\begin{answer}
D
\end{answer}


\begin{problem}
函数 \(f(x)=\cos2x+6\cos\left(\dfrac{\pi}{2}-x\right)\) 的最大值为（\quad）

\choices
    {\(4\)}
    {\(5\)}
    {\(6\)}
    {\(7\)}
\end{problem}

\begin{answer}
B
\end{answer}


\begin{problem}
已知函数 \(f(x)\)（\(x\in\R\)）满足 \(f(x)=f(2-x)\)，若函数 \(y=\lvert x^2-2x-3\rvert\) 与 \(y=f(x)\) 图象的交点为 \((x_1,y_1)\)，\((x_2,y_2)\)，\(\dots\)，\((x_m,y_m)\)，则 \(\displaystyle\sum_{i=1}^m x_i=\)（\quad）

\choices
    {\(0\)}
    {\(m\)}
    {\(2m\)}
    {\(4m\)}
\end{problem}

\begin{answer}
B
\end{answer}


\section{填空题}

\begin{problem}
已知向量 \(\bs a=(m,4)\)，\(\bs b=(3,-2)\)，且 \(\bs a\parallel\bs b\)，则 \(m=\fillinblank{}\).
\end{problem}

\begin{answer}
\(-6\)
\end{answer}


\begin{problem}
若 \(x\)，\(y\) 满足约束条件
\(\begin{cases}
x-y+1\ge 0\\
x+y-3\ge 0\\
x-3\le 0
\end{cases}\)
则 \(z=x-2y\) 的最小值为\fillinblank{}.
\end{problem}

\begin{answer}
\(-5\)
\end{answer}


\begin{problem}
\(\triangle ABC\) 的内角 \(A\)，\(B\)，\(C\) 的对边分别为 \(a\)，\(b\)，\(c\)，若 \(\cos A=\dfrac45\)，\(\cos C=\dfrac{5}{13}\)，\(a=1\)，则 \(b=\fillinblank{}\).
\end{problem}

\begin{answer}
\(\dfrac{21}{13}\)
\end{answer}


\begin{problem}
有三张卡片，分别写有 \(1\) 和 \(2\)，\(1\) 和 \(3\)，\(2\) 和 \(3\). 甲，乙，丙三人各取走一张卡片，甲看了乙的卡片后说：“我与乙的卡片上相同的数字不是 \(2\)”；乙看了丙的卡片后说：“我与丙的卡片上相同的数字不是 \(1\)”；丙说：“我的卡片上的数字之和不是 \(5\)”，则甲的卡片上的数字是\fillinblank{}.
\end{problem}

\begin{answer}
\(1\) 和 \(3\)
\end{answer}


\section{解答题}

\begin{problem}
等差数列 \(\{a_n\}\) 中，\(a_3+a_4=4\)，\(a_5+a_7=6\).

\begin{enumerate}
\item 求 \(\{a_n\}\) 的通项公式；
\item 设 \(b_n=[a_n]\)，求数列 \(\{b_n\}\) 的前 \(10\) 项和，其中 \([x]\) 表示不超过 \(x\) 的最大整数，如 \([0.9]=0\)，\([2.6]=2\).
\end{enumerate}
\end{problem}

\begin{solution}
\begin{enumerate}
\item 设等差数列 \(\{a_n\}\) 的首项为 \(a_1\)，公差为 \(d\)，则
\[
a_3+a_4=2a_1+5d=4.
\]
同理，
\[
a_5+a_7=2a_1+10d=6.
\]
两式相减得 \(5d=6-4=2\)，所以 \(d=\dfrac25\). 代入第一式，得
\[
2a_1+5\cdot\dfrac25=4\Longrightarrow a_1=1.
\]
因此
\[
a_n=a_1+(n-1)d=1+\dfrac25(n-1)=\dfrac{2n+3}{5}.
\]

\item 由 \(b_n=[a_n]\) 及上式，当 \(n=1,2,3\) 时有 \(1\leq a_n<2\)，故对应的 \(b_n\) 均为 \(1\)；当 \(n=4,5\) 时有 \(2\leq a_n<3\)，故对应的 \(b_n\) 均为 \(2\)；当 \(n=6,7,8\) 时有 \(3\leq a_n<4\)，故对应的 \(b_n\) 均为 \(3\)；当 \(n=9,10\) 时有 \(4\leq a_n<5\)，故对应的 \(b_n\) 均为 \(4\). 因而
\[
(b_1,b_2,\ldots,b_{10})=(1,1,1,2,2,3,3,3,4,4).
\]
所以数列 \(\{b_n\}\) 的前 \(10\) 项和为
\[
\sum_{n=1}^{10}b_n=3\cdot1+2\cdot2+3\cdot3+2\cdot4=24.
\]
\end{enumerate}
\end{solution}


\begin{problem}
某险种的基本保费为 \(a\)（单位：元），继续购买该险种的投保人称为续保人，续保人本年度的保费与其上年度出险次数的关联如下：

\begin{center}
\begin{tabular}{c|cccccc}
上年度出险次数 & \(0\) & \(1\) & \(2\) & \(3\) & \(4\) & \(\geq 5\)\\
\hline
保费 & \(0.85a\) & \(a\) & \(1.25a\) & \(1.5a\) & \(1.75a\) & \(2a\)
\end{tabular}
\end{center}

随机调查了该险种的 \(200\) 名续保人在一年内的出险情况，得到如下统计表：

\begin{center}
\begin{tabular}{c|cccccc}
上年度出险次数 & \(0\) & \(1\) & \(2\) & \(3\) & \(4\) & \(\geq 5\)\\
\hline
频数 & \(60\) & \(50\) & \(30\) & \(30\) & \(20\) & \(10\)
\end{tabular}
\end{center}

\begin{enumerate}
\item 记 \(A\) 为事件：“一续保人本年度的保费不高于基本保费”，求 \(P(A)\) 的估计值；
\item 记 \(B\) 为事件：“一续保人本年度的保费高于基本保费但不高于基本保费的 \(160\%\)”，求 \(P(B)\) 的估计值；
\item 求续保人本年度的平均保费估计值.
\end{enumerate}
\end{problem}

\begin{solution}
\begin{enumerate}
\item 由于 \(0.85a\leq a\)，而出险次数为 \(2\) 或更多时保费倍数均大于 \(1\)，本年度保费不高于基本保费的续保人对应上年度出险次数为 \(0\) 或 \(1\). 因此 \(P(A)\) 的估计值为
\[
P(A)\approx\dfrac{60+50}{200}=0.55.
\]

\item 本年度保费高于基本保费但不高于基本保费的 \(160\%\) 的续保人对应上年度出险次数为 \(2\) 或 \(3\)，因为此时保费倍数分别为 \(1.25\) 和 \(1.5\)，而出险次数为 \(4\) 或更多时保费倍数至少为 \(1.75\). 因此 \(P(B)\) 的估计值为
\[
P(B)\approx\dfrac{30+30}{200}=0.30.
\]

\item 以样本中各类续保人的保费加权平均数作为续保人本年度平均保费的估计值，得
\[
\begin{aligned}
\overline p&\approx\dfrac{60\cdot0.85a+50\cdot a+30\cdot1.25a+30\cdot1.5a+20\cdot1.75a+10\cdot2a}{200}\\
&=\dfrac{(51+50+37.5+45+35+20)a}{200}\\
&=\dfrac{238.5a}{200}=\dfrac{477}{400}a.
\end{aligned}
\]
\end{enumerate}
\end{solution}


\begin{problem}
如图，菱形 \(ABCD\) 的对角线 \(AC\) 与 \(BD\) 交于点 \(O\)，点 \(E\)，\(F\) 分别在 \(AD\)，\(CD\) 上，\(AE=CF\)，\(EF\) 交 \(BD\) 于点 \(H\)，将 \(\triangle DEF\) 沿 \(EF\) 折到 \(\triangle D'EF\) 的位置.
\begin{enumerate}
\item 证明：\(AC\perp HD'\)；
\item 若 \(AB=5\)，\(AC=6\)，\(AE=\dfrac54\)，\(OD'=2\sqrt2\)，求五棱锥 \(D'-ABCFE\) 的体积.
\end{enumerate}

\begingroup
\bitmapfigure[max width=0.34\linewidth]{img/2016/national_paper_2_liberal/q19_fig1.png}
\endgroup
\end{problem}

\begin{solution}
\begin{enumerate}
\item 因为 \(ABCD\) 是菱形，所以 \(AD=CD\). 又因为 \(AE=CF\)，所以
\[
\dfrac{AE}{AD}=\dfrac{CF}{CD}.
\]
在三角形 \(ADC\) 中，点 \(E\)，\(F\) 分别在边 \(AD\)，\(CD\) 上，由比例线段的逆定理可得 \(EF\parallel AC\). 菱形的对角线互相垂直，故 \(BD\perp AC\). 又因为 \(H\) 在 \(BD\) 和 \(EF\) 上，所以折叠前 \(DH\perp EF\). 折叠是绕直线 \(EF\) 的旋转，点 \(H\) 在折痕上保持不动，且线段与旋转轴的夹角不变，因此折叠后仍有 \(HD'\perp EF\). 因为 \(EF\parallel AC\)，所以
\[
AC\perp HD'.
\]

\item 在原平面内以 \(AC\) 和 \(BD\) 所在直线为坐标轴，取 \(O\) 为原点. 由于 \(OA=\dfrac{AC}{2}=3\)，
取 \(A=(-3,0)\)，\(C=(3,0)\).
在直角三角形 \(AOB\) 中，由勾股定理得
\[
OB=\sqrt{AB^2-OA^2}=\sqrt{25-9}=4.
\]
取 \(B=(0,4)\)，\(D=(0,-4)\).
于是 \(BD=8\). 由于 \(AD=5\)，且 \(AE=\dfrac54=\dfrac14AD\)，同理 \(CF=\dfrac14CD\)，所以
\[
E=\left(-\dfrac94,-1\right).
\]
同理
\[
F=\left(\dfrac94,-1\right).
\]
直线 \(EF\) 为 \(y=-1\)，直线 \(BD\) 为 \(x=0\)，所以
\[
H=(0,-1).
\]
从而 \(EF=\dfrac92\)，且点 \(D\) 到直线 \(EF\) 的距离为 \(3\). 折叠保持长度，故 \(HD'=HD=3\). 在三维直角坐标系中令原平面为 \(z=0\). 由第（1）问知 \(HD'\perp AC\)，而 \(AC\) 是 \(x\) 轴，故可设
\[
D'=(0,u,v).
\]
其中 \(\lvert v\rvert\) 是五棱锥的高. 由 \(OD'=2\sqrt2\) 和 \(HD'=3\)，分别有
\[
u^2+v^2=8.
\]
\[
(u+1)^2+v^2=9.
\]
两式相减得
\[
(u+1)^2-u^2=9-8\Longrightarrow 2u+1=1\Longrightarrow u=0.
\]
于是 \(v^2=8\)，所以五棱锥的高为 \(\lvert v\rvert=2\sqrt2\).

菱形 \(ABCD\) 的面积为
\[
S_{ABCD}=\dfrac12AC\cdot BD=\dfrac12\cdot6\cdot8=24.
\]
三角形 \(DEF\) 的底边为 \(EF=\dfrac92\)，高为点 \(D\) 到直线 \(EF\) 的距离 \(3\)，所以
\[
S_{DEF}=\dfrac12\cdot\dfrac92\cdot3=\dfrac{27}{4}.
\]
五边形 \(ABCFE\) 的面积为菱形面积减去三角形 \(DEF\) 的面积，即
\[
S_{ABCFE}=24-\dfrac{27}{4}=\dfrac{69}{4}.
\]
所以五棱锥 \(D'-ABCFE\) 的体积为
\[
V=\dfrac13S_{ABCFE}\cdot2\sqrt2=\dfrac{23\sqrt2}{2}.
\]
\end{enumerate}
\end{solution}


\begin{problem}
已知函数 \(f(x)=(x+1)\ln x-a(x-1)\).
\begin{enumerate}
\item 当 \(a=4\) 时，求曲线 \(y=f(x)\) 在 \((1,f(1))\) 处的切线方程；
\item 若当 \(x\in(1,+\infty)\) 时，\(f(x)>0\)，求 \(a\) 的取值范围.
\end{enumerate}
\end{problem}

\begin{solution}
\begin{enumerate}
\item 函数的定义域为 \(x>0\)，且
\[
f(1)=0.
\]
\[
f'(x)=\ln x+1+\dfrac1x-a.
\]
当 \(a=4\) 时，
\[
f'(1)=\ln1+1+1-4=-2.
\]
因此曲线在 \((1,f(1))=(1,0)\) 处的切线斜率为 \(-2\)，切线方程为
\[
y=-2(x-1).
\]
即 \(y=-2x+2\).

\item 当 \(x>1\) 时，令
\[
h(x)=(x+1)\ln x-2(x-1).
\]
则 \(h(1)=0\)，并且
\[
h'(x)=\ln x-1+\dfrac1x.
\]
\[
h''(x)=\dfrac{x-1}{x^2}>0.
\]
所以 \(h'(x)\) 在 \((1,+\infty)\) 上递增. 又因为 \(h'(1)=0\)，故当 \(x>1\) 时 \(h'(x)>0\)，进而 \(h(x)>h(1)=0\). 因此
\[
(x+1)\ln x>2(x-1).
\]
若 \(a\leq2\)，则对任意 \(x>1\)，有
\[
f(x)=h(x)+(2-a)(x-1)>0.
\]
反之，若 \(a>2\)，由洛必达法则得
\[
\lim_{x\to1^+}\dfrac{(x+1)\ln x}{x-1}=\lim_{x\to1^+}\left(\ln x+1+\dfrac1x\right)=2.
\]
从而
\[
\lim_{x\to1^+}\dfrac{f(x)}{x-1}=2-a<0.
\]
由于 \(x-1>0\)，所以在充分接近 \(1\) 的某些 \(x>1\) 处有 \(f(x)<0\)，不满足题设条件. 故 \(a\) 的取值范围为
\[
a\leq2.
\]
\end{enumerate}
\end{solution}


\begin{problem}
已知椭圆 \(E:\dfrac{x^2}{4}+\dfrac{y^2}{3}=1\)，\(A\) 是 \(E\) 的左顶点，斜率为 \(k\)（\(k>0\)）的直线交 \(E\) 于 \(A\)，\(M\) 两点，点 \(N\) 在 \(E\) 上，且 \(MA\perp NA\).
\begin{enumerate}
\item 当 \(\lvert AM\rvert=\lvert AN\rvert\) 时，求 \(\triangle AMN\) 的面积；
\item 当 \(2\lvert AM\rvert=\lvert AN\rvert\) 时，证明：\(\sqrt3<k<2\).
\end{enumerate}
\end{problem}

\begin{solution}
设 \(A=(-2,0)\). 过 \(A\) 的斜率为 \(k\) 的直线方向向量为 \((1,k)\)，故可设
\[
M=(-2+u,ku).
\]
将其代入椭圆方程，得
\[
\dfrac{(-2+u)^2}{4}+\dfrac{k^2u^2}{3}=1
\Longleftrightarrow 1-u+\dfrac{u^2}{4}+\dfrac{k^2u^2}{3}=1
\]
即 \(u\left(\dfrac{3+4k^2}{12}u-1\right)=0\).
由于 \(M\neq A\)，所以 \(u\neq0\)，从而
\[
u=\dfrac{12}{3+4k^2}.
\]
又因为 \(MA\perp NA\)，可取垂直于 \((1,k)\) 的方向向量 \((-k,1)\)，设
\[
N=(-2-kv,v).
\]
代入椭圆方程，得
\[
\dfrac{(-2-kv)^2}{4}+\dfrac{v^2}{3}=1
\Longleftrightarrow 1+kv+\dfrac{k^2v^2}{4}+\dfrac{v^2}{3}=1.
\]
即 \(v\left(k+\dfrac{3k^2+4}{12}v\right)=0\).
由于 \(N\neq A\)，所以 \(v\neq0\)，从而
\[
v=-\dfrac{12k}{3k^2+4}.
\]
因此
\[
|AM|=\dfrac{12\sqrt{1+k^2}}{3+4k^2}.
\]
\[
|AN|=\dfrac{12k\sqrt{1+k^2}}{3k^2+4}.
\]
\begin{enumerate}
\item 当 \(|AM|=|AN|\) 时，因 \(k>0\)，有
\[
\dfrac1{3+4k^2}=\dfrac{k}{3k^2+4}
\Longrightarrow 4k^3-3k^2+3k-4=0.
\]
因式分解得
\[
(k-1)(4k^2+k+4)=0.
\]
而 \(4k^2+k+4>0\)，所以 \(k=1\). 此时
\[
|AM|=\dfrac{12\sqrt2}{7}.
\]
由于 \(MA\perp NA\) 且 \(|AM|=|AN|\)，所以
\[
S_{\triangle AMN}=\dfrac12|AM|\cdot|AN|=\dfrac12\left(\dfrac{12\sqrt2}{7}\right)^2=\dfrac{144}{49}.
\]

\item 当 \(2|AM|=|AN|\) 时，得
\[
\dfrac2{3+4k^2}=\dfrac{k}{3k^2+4}
\Longrightarrow 4k^3-6k^2+3k-8=0.
\]
令 \(F(k)=4k^3-6k^2+3k-8\)，则
\[
F'(k)=12k^2-12k+3=3(2k-1)^2\geq0.
\]
且除 \(k=\dfrac12\) 外均有 \(F'(k)>0\)，所以 \(F(k)\) 在 \((0,+\infty)\) 上严格递增. 又
\[
F(\sqrt3)=15\sqrt3-26<0.
\]
并且
\[
F(2)=6>0.
\]
其中 \(15\sqrt3<26\) 因为 \((15\sqrt3)^2=675<676=26^2\). 由零点存在性及严格单调性可得
\[
\sqrt3<k<2.
\]
\end{enumerate}
\end{solution}


\section{选考题：解答题}

\begin{problem}
如图，在正方形 \(ABCD\) 中，\(E\)，\(G\) 分别在边 \(DA\)，\(DC\) 上（不与端点重合），且 \(DE=DG\)，过 \(D\) 点作 \(DF\perp CE\)，垂足为 \(F\).
\begin{enumerate}
\item 证明：\(B\)，\(C\)，\(G\)，\(F\) 四点共圆；
\item 若 \(AB=1\)，\(E\) 为 \(DA\) 的中点，求四边形 \(BCGF\) 的面积.
\end{enumerate}

\begingroup
\bitmapfigure[max width=0.34\linewidth]{img/2016/national_paper_2_liberal/q22_fig1.png}
\endgroup
\end{problem}

\begin{solution}
\begin{enumerate}
\item 设正方形边长为 \(s\)，并设 \(DE=DG=d\)，则 \(0<d<s\). 建立直角坐标系，取
\[
A=(0,0).
\]
取 \(B=(s,0)\)，\(C=(s,s)\)，\(D=(0,s)\). 于是
\[
E=(0,s-d).
\]
又有 \(G=(d,s)\).
设 \(F=E+\lambda(C-E)\)，则
\[
F=\bigl(\lambda s,s-d+\lambda d\bigr).
\]
因为 \(DF\perp CE\)，有
\[
\bigl(-\lambda s,d(1-\lambda)\bigr)\cdot(s,d)=0
\]
从而
\[
\lambda=\dfrac{d^2}{s^2+d^2}.
\]
设过 \(B\)，\(C\)，\(G\) 的圆方程为
\[
x^2+y^2+px+qy+r=0.
\]
将 \(B\)，\(C\) 的坐标分别代入，作差得
\[
\bigl(2s^2+ps+qs+r\bigr)-\bigl(s^2+ps+r\bigr)=0
\Longrightarrow q=-s.
\]
再将 \(B\)，\(G\) 的坐标代入，得
\[
s^2+ps+r=0.
\]
\[
d^2+pd+r=0.
\]
两式相减，且因 \(d\neq s\)，得 \(p=-(s+d)\)，进而由第一式得 \(r=sd\). 因此圆方程为
\[
\Phi(x,y)=x^2+y^2-(s+d)x-sy+sd=0.
\]
将 \(F\) 的坐标代入，得
\[
\Phi(F)=d^2-\lambda(s^2+2d^2)+\lambda^2(s^2+d^2).
\]
由 \(\lambda(s^2+d^2)=d^2\)，上式可化为
\[
\Phi(F)=(1-\lambda)\bigl(d^2(1-\lambda)-\lambda s^2\bigr)=0.
\]
所以 \(F\) 在过 \(B\)，\(C\)，\(G\) 的圆上，故 \(B\)，\(C\)，\(G\)，\(F\) 四点共圆.

\item 当 \(AB=1\) 且 \(E\) 为 \(DA\) 的中点时，\(s=1\)，\(d=DE=\dfrac12\). 因而
\[
\lambda=\dfrac{(1/2)^2}{1+(1/2)^2}=\dfrac15.
\]
由 \(F=\bigl(\lambda s,s-d+\lambda d\bigr)\)，得
\[
F=\left(\dfrac15,\dfrac35\right).
\]
同时
\[
B=(1,0).
\]
\[
C=(1,1).
\]
\[
G=\left(\dfrac12,1\right).
\]
按顶点顺序 \(B\)，\(C\)，\(G\)，\(F\) 使用鞋带公式，得
\[
S_{BCGF}=\dfrac12\left|\left(1+1+\dfrac3{10}\right)-\left(\dfrac12+\dfrac15+\dfrac35\right)\right|=\dfrac12.
\]
\end{enumerate}
\end{solution}


\begin{problem}
在直角坐标系 \(xOy\) 中，圆 \(C\) 的方程为 \((x+6)^2+y^2=25\).
\begin{enumerate}
\item 以坐标原点为极点，\(x\) 轴正半轴为极轴建立极坐标系，求 \(C\) 的极坐标方程；
\item 直线 \(l\) 的参数方程是
\(\begin{cases}
x=t\cos\alpha\\
y=t\sin\alpha
\end{cases}\)
（\(t\) 为参数），\(l\) 与 \(C\) 交于 \(A\)，\(B\) 两点，\(\lvert AB\rvert=\sqrt{10}\)，求 \(l\) 的斜率.
\end{enumerate}
\end{problem}

\begin{solution}
\begin{enumerate}
\item 令 \(x=\rho\cos\theta\)，\(y=\rho\sin\theta\)，代入圆的方程得
\[
(\rho\cos\theta+6)^2+\rho^2\sin^2\theta=25.
\]
整理得圆 \(C\) 的极坐标方程为
\[
\rho^2+12\rho\cos\theta+11=0.
\]

\item 直线 \(l\) 上的点满足 \(x=t\cos\alpha\)，\(y=t\sin\alpha\)，代入圆的方程得交点参数 \(t_1,t_2\) 满足
\[
t^2+12t\cos\alpha+11=0.
\]
由于 \((\cos\alpha,\sin\alpha)\) 是单位向量，弦长为 \(\lvert t_1-t_2\rvert\). 由根与系数的关系，
\[
t_1+t_2=-12\cos\alpha.
\]
\[
t_1t_2=11.
\]
因此
\[
\lvert AB\rvert^2=(t_1-t_2)^2=(t_1+t_2)^2-4t_1t_2=144\cos^2\alpha-44.
\]
由 \(\lvert AB\rvert=\sqrt{10}\)，得
\[
144\cos^2\alpha-44=10\Longrightarrow \cos^2\alpha=\dfrac38.
\]
所以 \(\sin^2\alpha=\dfrac58\)，从而直线 \(l\) 的斜率满足
\[
\tan^2\alpha=\dfrac{\sin^2\alpha}{\cos^2\alpha}=\dfrac53.
\]
因此 \(l\) 的斜率为
\[
\pm\sqrt{\dfrac53}=\pm\dfrac{\sqrt{15}}{3}.
\]
\end{enumerate}
\end{solution}


\begin{problem}
已知函数 \(f(x)=\lvert x-\dfrac12\rvert+\lvert x+\dfrac12\rvert\)，\(M\) 为不等式 \(f(x)<2\) 的解集.
\begin{enumerate}
\item 求 \(M\)；
\item 证明：当 \(a\)，\(b\in M\) 时，\(\lvert a+b\rvert<\lvert 1+ab\rvert\).
\end{enumerate}
\end{problem}

\begin{solution}
\begin{enumerate}
\item 由绝对值的定义，
\[
f(x)=\begin{cases}
-2x,&x\leq-\dfrac12,\\
1,&-\dfrac12\leq x\leq\dfrac12,\\
2x,&x\geq\dfrac12.
\end{cases}
\]
当 \(x\leq-\dfrac12\) 时，\(f(x)<2\) 等价于 \(-2x<2\)，即 \(x>-1\)，故这一段的解集为 \((-1,-\dfrac12]\). 当 \(-\dfrac12\leq x\leq\dfrac12\) 时，\(f(x)=1<2\)，这一段全部满足不等式. 当 \(x\geq\dfrac12\) 时，\(f(x)<2\) 等价于 \(2x<2\)，即 \(x<1\)，故这一段的解集为 \([\dfrac12,1)\). 合并三段得
\[
M=(-1,1).
\]

\item 当 \(a,b\in M\) 时，\(-1<a<1\)，\(-1<b<1\)，所以 \(1-a^2>0\) 且 \(1-b^2>0\)，从而
\[
(1-a^2)(1-b^2)>0.
\]
展开并移项，得
\[
(1+ab)^2-(a+b)^2>0.
\]
又因为 \(-1<ab<1\)，所以 \(1+ab>0\)，即 \(\lvert1+ab\rvert=1+ab\). 于是
\[
\lvert a+b\rvert^2<(1+ab)^2=\lvert1+ab\rvert^2.
\]
两边均非负，故 \(\lvert a+b\rvert<\lvert1+ab\rvert\).
\end{enumerate}
\end{solution}
